\section{Results}
\label{sec:nmi-results}
% Alias label so the reused appendix (SI) cross-references resolve in this variant.
\label{sec:results}

We grade the submitted PDF only; the model never sees OpenReview reviews, decisions,
ratings, venue tags, or author identities (Fig.~\ref{fig:design}). The two pre-registered
questions are not on equal footing. The first is the headline: does the AIPR score agree
with the human outcome? That agreement licenses the triage claim regardless of the second
question, which asks what the engineered pipeline adds over the model alone. We lead on the
large cohort $M$ ($n=\NminiPrimary$, AIPR on GPT-5.4-mini) and confirm on the nested cohort
$H$ ($n=\NfullPrimary$, AIPR on the production GPT-5.4); Table~\ref{tab:headline} gives both.

\begin{figure}[t]
  \centering
  \includegraphics[width=\textwidth]{fig0_design}
  \caption{\textbf{Study design and model-access boundary.} Each OpenReview submission
  supplies both the artifact reviewers saw (the submitted PDF) and the ground truth
  (decision tier and mean reviewer rating). AIPR grades the PDF with the frozen v6 two-pass
  pipeline under two nested configurations---AIPR (GPT-5.4-mini) and the production AIPR
  (GPT-5.4)---plus a Direct one-paragraph baseline on GPT-5.4; outcomes enter only after
  scoring. No model is tuned on the venue, decisions, or ratings.}
  \label{fig:design}
\end{figure}

\begin{table}[t]
  \centering
  \caption{Headline metrics on ICLR~2026, for the power configuration AIPR (GPT-5.4-mini)
  and the production AIPR (GPT-5.4). Brackets are 95\% bootstrap confidence intervals.}
  \label{tab:headline}
  \input{tables/tab_headline}
\end{table}

\subsection{The score agrees with the human outcome}
\label{sec:res-validity}% alias for reused SI cross-references
\paragraph{Weak submissions are flagged.}
We report the deployable flag on the production model directly. Its lowest score quintile
(the strict band shrinks to $n=\bottomNFrontier$ after integer-score ties) is rejected at
\bottomRejectRateFrontier\% (95\% CI \bottomRejectCIFrontier), a
\bottomLiftFrontier$\times$ lift over the \baseRejectRate\% base rate, with no oral papers in
the band; the pattern holds at scale on cohort $M$ (\bottomRejectRate\%, lift
\bottomLift$\times$, 95\% CI \bottomLiftCI{} clearing the pre-registered rule). We compute the
flag on the production score because the cheap-to-frontier bridge is weakest exactly at the
low end where the flag fires, so grading it directly sidesteps the proxy. The reliability
curve falls monotonically as the score rises (Fig.~\ref{fig:validation}b), and at the venue's
natural prevalence (\natAcceptRate\% accept) the bottom-quintile reject precision is
\natBottomPrecision\%.

\paragraph{Separation, gradient, and rating agreement.}
The overall score separates reject from accept at AUROC~\aurocFullFull{} on the production
cohort $H$, and \aurocMiniFull{} at scale on cohort $M$ (ROC in Fig.~\ref{fig:naive}a); the
accepted--rejected gap is large (Cohen's $d=\cohensDMini$, Cliff's
$\delta=\cliffsDeltaMini$). Mean score rises monotonically across the three tiers
(Jonckheere--Terpstra $\trendPrel$; Spearman $\rho=\trendRho$, 95\% CI \trendRhoCI), and it
tracks the continuous mean reviewer rating at $\rho=$~\spearmanRatingMiniFull{} (cohort $M$)
and \spearmanRatingFullFull{} (cohort $H$). Because the rating averages over reviewer noise,
this is the cleaner test of agreement with human-perceived quality, and the effect is strong
rather than merely significant (Fig.~\ref{fig:validation}). The agreement shows plainly at the
score range's ends: AIPR's top-scoring submission was an accepted oral (overall~83), its
lowest a reject (overall~58, rating~3.5).

\begin{figure}[t]
  \centering
  \includegraphics[width=\textwidth]{fig_validation}
  \caption{\textbf{The AIPR score agrees with the human outcome three ways} (cohort $M$).
  \textbf{(a)} Overall score across the ordered decision tiers, rising monotonically.
  \textbf{(b)} Reject rate by score decile with Wilson 95\% intervals; shaded band = bottom
  score quintile, dashed line = base reject rate. \textbf{(c)} Score vs.\ mean reviewer
  rating, coloured by tier, with a least-squares trend and Spearman $\rho$.}
  \label{fig:validation}
\end{figure}

\paragraph{The cheap-model proxy is valid.}
On cohort $H$ the cheap and production scores correlate at $\rho=$~\bridgeRhoFull{}
($n=\bridgeN$), clearing the pre-registered $\rho\ge0.8$ gate. Because a high global
correlation can hide low-end disagreement, we check directly: within the cohort-$M$ bottom
quintile the two still correlate at $\rho=$~\bridgeLowRhoFull{}, and \bridgeOverlap\% of that
band is also in the production bottom quintile. Cohort $M$ therefore stands in for the
production score at scale, with cohort $H$ as confirmation.

\subsection{What the engineering adds: a strong model, made reliable}
\label{sec:res-value}
A frontier model is available to anyone, so what does the pipeline add over asking it
directly? The pre-registered value comparison grades cohort $H$ a second way, Direct
(GPT-5.4): the same model and PDF with one paragraph, no rubric, audit, or examples, so the
baseline is realistic rather than a strawman (verbatim prompt in the Supplementary
Information).

Direct (GPT-5.4) separates reject from accept on its own at AUROC~\aurocNaiveFull{}. AIPR
reaches \aurocFullFull{} on the same papers; the paired difference is small and favours the
pipeline but does not meet the pre-declared superiority criterion ($\Delta$AUROC
\aurocDiffFullNaive, 95\% CI \aurocDiffCI, \aurocDiffPrel; $n=\naiveN$), and a post-hoc power
analysis puts the smallest gap detectable at this $n$ above the observed one (SI). We claim
neither superiority nor equivalence. The validated agreement above is therefore not
manufactured by the engineering: with or without the rubric, a frontier model already
produces a first-pass score that tracks the human outcome, and raw intelligence is not the
bottleneck.

What the engineering buys is not a higher score but a stable one. Across repeated gradings of
the same paper, AIPR's overall barely moves---median within-paper SD \fullRunSD{} points,
against \naiveRunSD{} for Direct on the same model and papers (exact Wilcoxon signed-rank
\relPairedP{}, $n=\relPairedN$ pairs; Fig.~\ref{fig:naive}b). Where the discrimination gap is
unresolved, this reliability gap is large and unambiguous: the same submission earns nearly
the same verdict every time. And the same pass returns the structured review
itself---rubric-anchored dimensions, a grounded literature audit, verbatim-anchored
comments---which is the system's actual output; the score is a by-product, and Direct returns
only the number. That review, not the scalar, is what a reviewer acts on.

\begin{figure}[t]
  \centering
  \includegraphics[width=\textwidth]{fig_naive}
  \caption{\textbf{What the engineering adds: a strong model made reliable} (cohort $H$;
  pre-registered value comparison, same model and submitted PDF throughout). \textbf{(a)}
  Reject-vs-accept ROC for AIPR (GPT-5.4), AIPR (GPT-5.4-mini), and the Direct (GPT-5.4)
  baseline; the legend gives each grader's AUROC (95\% CIs in Table~\ref{tab:headline}), and
  the AIPR$-$Direct paired difference is not statistically resolved (\aurocDiffPrel), so the
  engineering is not buying raw discrimination. \textbf{(b)} Within-paper score SD across
  repeated runs, one dot per paper coloured by decision tier, with a median bar; AIPR grades
  far more consistently (median \fullRunSD{} vs.\ \naiveRunSD{} points), a gap that is large
  where the discrimination gap is not. Interactive counterpart:
  \url{https://aipr.pub/publications\#fig-naive}.}
  \label{fig:naive}
\end{figure}

\subsection{Secondary analyses and failure modes}
\label{sec:failure}% alias for reused SI cross-references
Per-dimension discrimination is comparable across the four informative dimensions, with rigor
and applicability highest and novelty lowest (SI); the lightly weighted citation dimension is
uninformative this run and excluded from interpretation, and the headline is unchanged when it
is dropped (AUROC \aurocDropCitationFull{} vs.\ \aurocFull{}). A triage signal must be
characterized where it fails. Reading the public reviews of cohort $H$ (full case studies in
the SI), the dangerous error for triage---a weak paper scored high---concentrates where
rejection rests on a contribution-level judgment a competent manuscript surface does not
betray, or on a defect the system cannot see (it cannot run an experiment or re-derive a
proof). The mirror error, strong work scored low, reflects a genuine concern a reviewer also
raised but the community forgave, and costs only a second human look. That asymmetry---cheap
false positives, expensive false negatives at the top---is what makes the narrow ``flag weak
work'' claim defensible, and why the validated use is surfacing weak work for attention, never
automated rejection.
